%!TEX root = ../main.tex
% =============================================================
%  CAPÍTULO 3 — MARCO TEÓRICO Y ESTADO DEL ARTE
%                                       (Entrega 3 · 10 % · 28/7)
% =============================================================

\chapter{Síntesis de la literatura consultada}\label{cap:marco-teorico}

\begin{guia}[Guía de la cátedra: qué se evalúa en este capítulo]
\begin{enumerate}
    \item \textbf{Estado del arte}: investigaciones nacionales,
          internacionales y productos similares, por separado. De cada
          trabajo: qué hizo y qué resultados obtuvo, con métricas.
    \item \textbf{Brechas identificadas}: qué falta en lo existente y
          justifica \emph{este} proyecto. Es lo que separa un estado del
          arte de un listado. Incluir también los resultados negativos de
          la literatura.
    \item \textbf{Bases teóricas}: desarrollan las variables del título y su
          relación. Cada concepto termina explicando cómo se usa en el
          proyecto (teoría instrumental, no de relleno). Citas canónicas
          antes que tutoriales.
    \item \textbf{Marco conceptual}: conceptos que se nombran pero no se
          desarrollan. Los términos del glosario se marcan con
          \verb|\gls{}|, por ejemplo \gls{g:mvp} y \gls{prototipo}.
    \item \textbf{Calidad bibliográfica}: papers, tesis y libros de
          editoriales serias, con datos completos. Wikipedia y blogs restan.
    \item \textbf{Entrevistas a especialistas} (si las hay): fecha, nombre,
          cargo y aporte; transcripción en el apéndice~\ref{apx:entrevista}.
\end{enumerate}
Errores que la cátedra marca: falta de coherencia interna, bibliografía
superficial, información irrelevante o redundante. Es uno de los
capítulos más extensos (orientativo: 20 a 30 páginas).
\end{guia}

\section{Seguridad vial y exposición de motociclistas}

Las estadísticas oficiales y los estudios académicos consultados
coinciden en un punto: por kilómetro recorrido, el motociclista es el
usuario vial más expuesto a sufrir una lesión grave. El Global status
report on road safety 2023, de la Organización Mundial de la Salud,
ubica a peatones, ciclistas y motociclistas como los grupos que
concentran de manera desproporcionada las muertes de tránsito a nivel
global, en especial en países de ingresos bajos y medios, y señala que
los siniestros viales siguen siendo la principal causa de muerte de
niños y jóvenes de 5 a 29 años (WHO, 2023).

Para el contexto latinoamericano, CAF (2013) propone una metodología
específica para que los gobiernos locales elaboren planes de seguridad
vial dirigidos a motociclistas, con un enfoque participativo que integra
diagnóstico, diseño de medidas y seguimiento. El documento parte de una
premisa que MotoSafe GBA comparte: la motocicleta no puede tratarse como
una variante menor del automóvil dentro de la política de seguridad
vial, porque su patrón de siniestralidad y de exposición es
cualitativamente distinto.

En Argentina, datos oficiales de la Agencia Nacional de Seguridad Vial
correspondientes a 2025 ubicaron a los motociclistas en el 46\% de las
víctimas fatales de siniestros viales del país (Infoecos, 2026, citando
a la ANSV), en línea con coberturas periodísticas independientes que
reportaron cifras de 46-47\% para el mismo período (La Nación, 2025;
Dato Posta, 2025). \citet{camacho2025}, en su tesis de maestría sobre
implementación del Plan Estratégico de Seguridad Vial 2024-2027 en
gobiernos locales argentinos (incluyendo casos del AMBA), documenta que
la capacidad institucional para ejecutar políticas de seguridad vial,
entre ellas las dirigidas a motociclistas, depende tanto del diseño
nacional como del desarrollo institucional y la articulación
intersectorial a nivel municipal. Este antecedente sitúa a MotoSafe GBA
como una herramienta de nivel operativo (información al usuario final)
que no sustituye, pero puede complementar, ese nivel de política
pública.

\section{Datos geoespaciales, cobertura y comunicación de incertidumbre}

Un estudio de \citet{alsahfi2024} analiza patrones de siniestralidad vial en
ciudades de California durante un período de cinco años combinando tres
herramientas de análisis espacial: DBSCAN (un algoritmo de clustering
que agrupa puntos por densidad, sin necesidad de fijar de antemano
cuántos grupos existen), estimación de densidad de kernel o KDE (que
suaviza puntos dispersos en un mapa de densidad continuo) y el
estadístico Getis-Ord Gi* (que identifica, con significancia
estadística, si una zona concentra más o menos casos de los esperables
por azar). El aporte metodológico relevante para este proyecto no es
solo la técnica de clustering, similar a la usada para los hotspots de
inseguridad de MotoSafe GBA, sino la insistencia del estudio en reportar
la relación entre densidad poblacional y siniestralidad detectada: sin
ese control, un área con más registros puede confundirse con un área más
riesgosa cuando en realidad solo tiene más tránsito o más reporte.

La comunicación de incertidumbre en datos de origen colaborativo tiene
un antecedente directo en la evaluación de reportes de Waze realizada
por \citet{aminnaseri2018}, que midió que los reportes colaborativos
de esa aplicación cubrían el 43,2\% de los eventos de choque y
congestión registrados por sistemas oficiales de gestión de tránsito. La
cifra es útil no como benchmark de cobertura a igualar, sino como
evidencia de que una fuente colaborativa aporta valor real sin alcanzar
cobertura total, y de que ese porcentaje debe comunicarse en lugar de
asumirse. MotoSafe GBA aplica el mismo principio a su capa de prensa y
de reportes comunitarios: la ausencia de un registro en una zona se
muestra como "sin datos", explícitamente distinto de un "0\%" que
sugeriría ausencia de riesgo.

\section{Aprendizaje automático para gravedad vial}

Un estudio de \citet{islam2022} compara tres modelos de clasificación
(Random Forest, Gradient Boosting y regresión logística) para predecir
la severidad de siniestros viales en Al-Ahsa, Arabia Saudita, sobre un
mismo conjunto de datos de siniestros con víctimas, e incorpora
autocorrelación espacial mediante el estadístico Getis-Ord G* para
localizar zonas de alta y baja severidad. En esa comparación directa, el
Random Forest obtuvo mejor desempeño de clasificación que la regresión
logística y que el Gradient Boosting, con el tipo de colisión como la
variable de mayor peso en sus predicciones; ese resultado respalda la
elección de un Random Forest como algoritmo base para MotoSafe GBA, en
vez de partir directamente de un modelo más simple (regresión logística)
o de uno más costoso de ajustar (Gradient Boosting) sin evidencia previa
de que la complejidad adicional se traduzca en mejor desempeño sobre
datos de este tipo. Ese estudio, sin embargo, no reporta el desbalance
de clases que presenta el dataset de este proyecto (LEVE muy mayoritaria
frente a GRAVE y MORTAL), lo que exige un tratamiento adicional durante
el entrenamiento y una lectura cuidadosa de precisión frente a recall
por clase.

La documentación oficial de scikit-learn (la biblioteca de Python de
referencia para aprendizaje automático clásico, la misma que se usa para
entrenar y evaluar el modelo de MotoSafe GBA) sobre errores frecuentes
de evaluación \citep{sklearndocs} formaliza un riesgo
metodológico central para este tipo de proyectos: la fuga de datos (data
leakage) que ocurre cuando una transformación (selección de variables,
imputación, codificación) se ajusta sobre todo el conjunto de datos
antes de separar entrenamiento y evaluación, inflando artificialmente el
desempeño reportado. El protocolo de evaluación temporal descripto en el
Capítulo 2 sigue esa recomendación de manera explícita: la imputación de
edad y la codificación de categorías se ajustan únicamente sobre el
conjunto de entrenamiento.

Random Forest es, además, una elección apropiada por el tipo de dato de
entrada del proyecto: MotoSafe GBA trabaja con datos tabulares (hora,
comuna, tipo de vía, clima, variables demográficas), no con imágenes ni
con secuencias temporales complejas. Un modelo de árboles como Random
Forest es más liviano de entrenar, no requiere infraestructura de GPU, y
permite auditar la importancia de cada variable por separado, algo
relevante para un proyecto que necesita poder explicar por qué el modelo
predice lo que predice, no solo qué predice.

\section{Sistemas de información geográfica, navegación y fuentes públicas}

Representar información vial requiere algo más que un mapa de fondo con
puntos dibujados encima: cada consulta necesita responderse contra
geometrías reales, idealmente con operaciones espaciales (distancia,
contención, intersección) resueltas por la propia base de datos en vez
de por el código de la aplicación. Un sistema de información geográfica
(SIG) es, en términos generales, cualquier sistema que almacena,
consulta y analiza datos asociados a una ubicación en el espacio; la
decisión de diseño relevante para este proyecto no es si usar un SIG,
sino en qué capa de la arquitectura vive esa lógica espacial: en la base
de datos, en el servidor de aplicación, o en el navegador del usuario.

MotoSafe GBA elige resolver esa lógica en la base de datos: se apoya en
PostGIS como extensión geoespacial de PostgreSQL, y en MapLibre GL JS
como motor de renderizado cartográfico sin dependencia de un proveedor
cerrado de mapas. Tratar el componente espacial como una capa de
análisis propia, con sus propias herramientas estadísticas, y no como un
simple fondo visual sobre el que se dibujan puntos, es lo que justifica
elegir PostGIS en primer lugar: sus funciones espaciales permiten
resolver proximidad y pertenencia a un hotspot directamente en la base,
en vez de recorrer listas de coordenadas en el backend. Hoy, con PostGIS
habilitado y dos tablas con geometría nativa, el cálculo de proximidad
corre en código de aplicación; migrarlo a funciones espaciales de
PostGIS queda como el siguiente paso natural de esa capa.

Los datos históricos de siniestros y de inseguridad describen un patrón
general de una zona, pero no dicen nada sobre lo que pasó ahí la semana
pasada: para eso, MotoSafe GBA suma una capa de noticias locales, que
aporta contexto reciente (un corte de calle, un operativo, un hecho
puntual) que los datasets históricos no capturan por su propia
naturaleza, ya que se actualizan con mucha menor frecuencia. Automatizar
la lectura periódica de varios sitios de noticias, en lugar de
revisarlos manualmente, exige respetar las reglas que cada sitio publica
sobre qué contenido puede recorrer un proceso automático y qué contenido
está reservado a la lectura humana.

\section{Estado del arte y brecha abordada por MotoSafe}

Los antecedentes revisados en este capítulo responden, cada uno, a una
pregunta distinta de la que responde MotoSafe GBA, y esa diferencia es
la que define el aporte del proyecto. Por un lado, los trabajos de
análisis espacial y de clasificación de severidad (Islam et al., 2022;
Alsahfi, 2024) están pensados para que un organismo de tránsito o un
planificador urbano identifique, a escala de ciudad, dónde concentrar
recursos o intervenciones: responden "¿qué zonas requieren atención?",
no "¿qué puedo saber yo, ahora, sobre esta zona antes de circular por
ahí?". Por otro lado, los documentos de política pública y planificación
institucional \citep{caf2013,camacho2025} responden una pregunta de
gestión: cómo diseñar e implementar un plan de seguridad vial a nivel
municipal o nacional, sin bajar nunca a una herramienta que un
motociclista concreto pueda consultar antes de un viaje.

Ninguno de los dos grupos de antecedentes construye, ni se propone
construir, una consulta orientada al usuario final que integre esas
fuentes con su procedencia y su cobertura declaradas. Esa es la brecha
que aborda MotoSafe GBA: no la ausencia total de datos, sino la
distancia entre fuentes dispersas y heterogéneas en procedencia y
cobertura, por un lado, y una consulta comprensible para un motociclista
concreto, por el otro. La evidencia sobre cobertura parcial de fuentes
colaborativas \citep{aminnaseri2018} y el riesgo metodológico de
fuga de datos en modelos de clasificación (scikit-learn developers,
s.f.) son, en ese sentido, restricciones de diseño tanto como
antecedentes teóricos: definen qué no puede afirmarse sin evidencia
adicional, y por eso delimitan tanto lo que el sistema muestra como lo
que se abstiene de afirmar.
